\documentclass[aspectratio=169]{beamer}
\usetheme{Madrid}
\usecolortheme{default}

\usepackage{amsmath,amssymb,amsfonts,bm,mathtools}
\usepackage{physics}
\usepackage{braket}

\title{Solving quantum mechanical problems with VQE algorithm}
\author{Morten Hjorth-Jensen}
\date{FYS5419/9419 Spring 2026, March 9-13}

\begin{document}

\frame{\titlepage}

\begin{frame}[plain,fragile]
\frametitle{Plans for the week of March 10-14: Solving quantum mechanical problems with VQE algorithm}

\begin{block}{}
\begin{enumerate}
\item Discussion of  our final model, the Lipkin model, a two-qubit and a four-qubit system, reminder from last week, code examples and theory



\item Jordan-Wigner transformation

\item Applications to other Hamiltonians, pairing Hamiltonian  (codes and theory)

\item If we get time: Start discussion of Fourier transforms, continuous and discrete,, general motivation

\item Discussions and work on project 1 during lab sessions
\item \href{{https://youtu.be/C8vxBY-AmD8}}{Video of lecture}
\item \href{{https://github.com/CompPhysics/QuantumComputingMachineLearning/blob/gh-pages/doc/HandWrittenNotes/2026/Notesweek8.pdf}}{Whiteboard notes}
\end{enumerate}

\noindent
\end{block}
\end{frame}



\begin{frame}{Hamiltonian in terms of Pauli strings}
We have seen during the previous weeks how to rewrite many-body Hamiltonians in terms of Pauli strings
built from
\[
I,\qquad X,\qquad Y,\qquad Z,
\]
so that they are ready for quantum-computing implementations.

\medskip

\begin{itemize}
\item the Lipkin Hamiltonian in quasispin form,
\item detailed mappings to Pauli operators,
\item compressed and direct qubit encodings,
\item second-quantized fermionic operators \(a^\dagger,a\),
\item the nuclear pairing Hamiltonian as a second worked example,
\item a fuller Jordan--Wigner mapping from fermions to qubits.
\end{itemize}
\end{frame}

\begin{frame}{Notation}
We use:
\begin{itemize}
\item \(a_p^\dagger\): fermion creation operator in orbital \(p\),
\item \(a_p\): fermion annihilation operator in orbital \(p\),
\item \(n_p=a_p^\dagger a_p\): number operator,
\item \(X_i,Y_i,Z_i\): Pauli operators on qubit \(i\),
\item \(I_i\): identity on qubit \(i\).
\end{itemize}

Throughout, we use only
\[
a_p^\dagger,\qquad a_p.
\]

For quasispin systems we use collective operators
\[
J_x,\; J_y,\; J_z,\qquad J_\pm = J_x \pm i J_y.
\]
For pairing we use local quasispin operators
\[
S_p^\pm,\qquad S_p^z.
\]
\end{frame}

\begin{frame}{Why Pauli decompositions matter}
Quantum algorithms typically require a Hamiltonian of the form
\[
H = \sum_\alpha h_\alpha P_\alpha,
\qquad
P_\alpha \in \{I,X,Y,Z\}^{\otimes n}.
\]

Once a Hamiltonian is written in this way, one can:
\begin{itemize}
\item estimate energies in VQE,
\item generate Trotterized time evolution,
\item build QAOA circuits,
\item exploit commutativity and measurement grouping.
\end{itemize}

So the main preprocessing step is:
\begin{quote}
rewrite the many-body Hamiltonian as a weighted sum of Pauli strings.
\end{quote}
\end{frame}

% =========================================================
\section{Part I: Lipkin Hamiltonian in quasispin form}
% =========================================================

\begin{frame}{Lipkin Hamiltonian}
We consider the Lipkin Hamiltonian in quasispin representation
\[
H
=
\frac{\epsilon}{2}\left(J_z-\frac{N}{2}\right)
+
\frac{V}{2}\left(J_+^2+J_-^2\right)
+
\frac{W}{2}J_z^2.
\]

Here:
\begin{itemize}
\item \(N\) is the number of particles,
\item \(\epsilon\) sets the one-body splitting,
\item \(V\) controls pair scattering between the two levels,
\item \(W\) weights the quadratic \(J_z^2\) term.
\end{itemize}
\end{frame}

\begin{frame}{Quasispin algebra}
The collective operators satisfy the \(su(2)\) algebra
\[
[J_x,J_y]=iJ_z,\qquad
[J_y,J_z]=iJ_x,\qquad
[J_z,J_x]=iJ_y,
\]
or, equivalently,
\[
[J_z,J_\pm]=\pm J_\pm,
\qquad
[J_+,J_-]=2J_z.
\]

The basis states are
\[
\ket{j,m},
\qquad
m=-j,-j+1,\dots,j,
\]
with
\[
J_z\ket{j,m}=m\ket{j,m}.
\]
\end{frame}

\begin{frame}{Action of ladder operators}
The ladder operators act as
\[
J_+\ket{j,m}
=
\sqrt{j(j+1)-m(m+1)}\ket{j,m+1},
\]
\[
J_-\ket{j,m}
=
\sqrt{j(j+1)-m(m-1)}\ket{j,m-1}.
\]

Therefore
\[
J_+^2\ket{j,m}
=
\sqrt{j(j+1)-m(m+1)}
\sqrt{j(j+1)-(m+1)(m+2)}
\ket{j,m+2},
\]
and analogously for \(J_-^2\).
\end{frame}

\begin{frame}{Useful identity for the interaction term}
Using
\[
J_\pm = J_x \pm iJ_y,
\]
we compute
\[
J_+^2 + J_-^2
=
(J_x+iJ_y)^2 + (J_x-iJ_y)^2
=
2(J_x^2-J_y^2).
\]

Hence the Lipkin Hamiltonian can also be written as
\[
H
=
\frac{\epsilon}{2}\left(J_z-\frac{N}{2}\right)
+
V(J_x^2-J_y^2)
+
\frac{W}{2}J_z^2.
\]
\end{frame}

\begin{frame}{Two mapping strategies}
There are two main strategies for turning the Lipkin Hamiltonian into Pauli strings:
\begin{enumerate}
\item \textbf{Compressed encoding of the quasispin basis}: encode the \((2j+1)\)-dimensional quasispin space into \(\lceil \log_2(2j+1)\rceil\) qubits.
\item \textbf{Direct collective-spin mapping}: write
\[
J_\alpha = \frac12\sum_i \sigma_i^\alpha
\]
on a multi-qubit register and expand the Hamiltonian.
\end{enumerate}

We will use:
\begin{itemize}
\item compressed encoding for \(N=4\),
\item direct collective mapping for \(N=6\).
\end{itemize}
\end{frame}

\subsection{Case \(N=4\): compressed 3-qubit encoding}

\begin{frame}{Four particles: quasispin dimension and encoding}
For \(N=4\),
\[
j=\frac{N}{2}=2,
\qquad
2j+1 = 5.
\]

Thus the quasispin sector has dimension \(5\), which fits into \(3\) qubits since
\[
2^2=4<5\le 8=2^3.
\]

We choose the binary map
\[
\ket{2,-2}\leftrightarrow \ket{000},\quad
\ket{2,-1}\leftrightarrow \ket{001},\quad
\ket{2,0}\leftrightarrow \ket{010},
\]
\[
\ket{2,1}\leftrightarrow \ket{011},\quad
\ket{2,2}\leftrightarrow \ket{100}.
\]

The states \(\ket{101},\ket{110},\ket{111}\) are unphysical.
\end{frame}

\begin{frame}{Projectors and transition operators on qubits}
For one qubit:
\[
\ket{0}\bra{0}=\frac{I+Z}{2},
\qquad
\ket{1}\bra{1}=\frac{I-Z}{2},
\]
\[
\ket{1}\bra{0}=\frac{X-iY}{2},
\qquad
\ket{0}\bra{1}=\frac{X+iY}{2}.
\]

Therefore:
\begin{itemize}
\item diagonal operators are expanded with \(I\) and \(Z\),
\item off-diagonal transitions are expanded with \(X\) and \(Y\),
\item any finite operator in the encoded space can be written as a sum of Pauli strings.
\end{itemize}
\end{frame}

\begin{frame}{Three-qubit form of \(J_z\)}
Since
\[
J_z = \sum_{m=-2}^2 m\ket{m}\bra{m},
\]
we obtain
\[
J_z
=
-2\ket{000}\bra{000}
-\ket{001}\bra{001}
+\ket{011}\bra{011}
+2\ket{100}\bra{100}.
\]

Expanding the projectors gives
\[
J_z
=
-\frac{1}{4} IZI
+\frac{1}{4} IZZ
-\frac{1}{2} ZII
-\frac{1}{2} ZIZ
-\frac{3}{4} ZZI
-\frac{1}{4} ZZZ.
\]
\end{frame}

\begin{frame}{Three-qubit form of \(J_z^2\)}
Similarly,
\[
J_z^2 = \sum_{m=-2}^2 m^2 \ket{m}\bra{m},
\]
which yields
\[
J_z^2
=
\frac{5}{4} III
+\frac{3}{4} IIZ
+IZI
+IZZ
+\frac{1}{4} ZII
-\frac{1}{4} ZIZ.
\]

These formulas are exact on the physical 5-dimensional subspace.
\end{frame}

\begin{frame}{Three-qubit form of the ladder terms}
For \(j=2\),
\[
J_+\ket{-2}=2\ket{-1},\qquad
J_+\ket{-1}=\sqrt{6}\ket{0},
\]
\[
J_+\ket{0}=\sqrt{6}\ket{1},\qquad
J_+\ket{1}=2\ket{2}.
\]

Therefore
\[
J_+^2+J_-^2
=
\frac{3+\sqrt{6}}{2}\,IXI
+\frac{\sqrt{6}-3}{2}\,IXZ
+\frac{\sqrt{6}}{2}\,XXI
+\frac{\sqrt{6}}{2}\,XXZ
\]
\[
\qquad
+\frac{\sqrt{6}}{2}\,YYI
+\frac{\sqrt{6}}{2}\,YYZ
+\frac{3+\sqrt{6}}{2}\,ZXI
+\frac{\sqrt{6}-3}{2}\,ZXZ.
\]
\end{frame}

\begin{frame}{Final \(N=4\), 3-qubit Lipkin Hamiltonian}
Substituting into
\[
H = \frac{\epsilon}{2}(J_z-2) + \frac{V}{2}(J_+^2+J_-^2) + \frac{W}{2}J_z^2
\]
gives
\[
H_{N=4}^{(3q)}
=
\left(-\epsilon+\frac{5W}{8}\right)III
+\frac{3W}{8}IIZ
+\left(-\frac{\epsilon}{8}+\frac{W}{2}\right)IZI
+\left(\frac{\epsilon}{8}+\frac{W}{2}\right)IZZ
\]
\[
+\left(-\frac{\epsilon}{4}+\frac{W}{8}\right)ZII
+\left(-\frac{\epsilon}{4}-\frac{W}{8}\right)ZIZ
-\frac{3\epsilon}{8}ZZI
-\frac{\epsilon}{8}ZZZ
\]
\[
+\frac{V(3+\sqrt{6})}{4}IXI
+\frac{V(\sqrt{6}-3)}{4}IXZ
+\frac{V\sqrt{6}}{4}XXI
+\frac{V\sqrt{6}}{4}XXZ
+\frac{V\sqrt{6}}{4}YYI
+\frac{V\sqrt{6}}{4}YYZ
+\frac{V(3+\sqrt{6})}{4}ZXI
+\frac{V(\sqrt{6}-3)}{4}ZXZ.
\]
\end{frame}

\begin{frame}{Code-space penalty}
Because the 3-qubit register has 8 basis states but only 5 are physical, one often adds
\[
H_{\mathrm{pen}} = \lambda P_{\mathrm{unphys}},
\]
where
\[
P_{\mathrm{unphys}}
=
\ket{101}\bra{101}
+\ket{110}\bra{110}
+\ket{111}\bra{111}.
\]

This energetically suppresses leakage outside the encoded quasispin space.
\end{frame}

\subsection{Case \(N=6\): direct 6-qubit mapping}

\begin{frame}{Six particles: direct collective-spin realization}
For \(N=6\), a simple and transparent mapping uses six qubits and defines
\[
J_x = \frac{1}{2}\sum_{i=1}^{6} X_i,
\qquad
J_y = \frac{1}{2}\sum_{i=1}^{6} Y_i,
\qquad
J_z = \frac{1}{2}\sum_{i=1}^{6} Z_i.
\]

This is not qubit-optimal, since the symmetric quasispin sector has dimension \(7\), but it leads to a very clean Pauli-string Hamiltonian.
\end{frame}

\begin{frame}{Expanding the collective squares}
Using \(X_i^2=Y_i^2=Z_i^2=I\),
\[
J_x^2
=
\frac{1}{4}\sum_{i=1}^{6} I
+
\frac{1}{4}\sum_{i\neq j} X_iX_j
=
\frac{6}{4}I + \frac{1}{2}\sum_{i<j} X_iX_j,
\]
and similarly
\[
J_y^2
=
\frac{6}{4}I + \frac{1}{2}\sum_{i<j} Y_iY_j,
\]
\[
J_z^2
=
\frac{6}{4}I + \frac{1}{2}\sum_{i<j} Z_iZ_j.
\]
\end{frame}

\begin{frame}{Interaction term on 6 qubits}
Using
\[
J_+^2+J_-^2=2(J_x^2-J_y^2),
\]
we get
\[
\frac{V}{2}(J_+^2+J_-^2)
=
V(J_x^2-J_y^2)
=
\frac{V}{2}\sum_{i<j}(X_iX_j-Y_iY_j).
\]

Thus the interaction is an all-to-all two-body Pauli interaction.
\end{frame}

\begin{frame}{Linear and quadratic \(J_z\) terms on 6 qubits}
For the linear term,
\[
\frac{\epsilon}{2}\left(J_z-\frac{N}{2}\right)
=
\frac{\epsilon}{2}\left(\frac12\sum_{i=1}^{6}Z_i - 3I\right)
=
\frac{\epsilon}{4}\sum_{i=1}^{6} Z_i - \frac{3\epsilon}{2}I.
\]

For the quadratic term,
\[
\frac{W}{2}J_z^2
=
\frac{W}{2}\left(\frac{6}{4}I+\frac12\sum_{i<j}Z_iZ_j\right)
=
\frac{3W}{4}I+\frac{W}{4}\sum_{i<j}Z_iZ_j.
\]
\end{frame}

\begin{frame}{Final \(N=6\), 6-qubit Lipkin Hamiltonian}
Collecting all terms gives
\[
H_{N=6}^{(6q)}
=
\left(-\frac{3\epsilon}{2}+\frac{3W}{4}\right)I
+
\frac{\epsilon}{4}\sum_{i=1}^{6} Z_i
+
\frac{V}{2}\sum_{i<j}(X_iX_j-Y_iY_j)
+
\frac{W}{4}\sum_{i<j} Z_iZ_j.
\]

This is already directly usable in qubit-based quantum algorithms.
\end{frame}

\begin{frame}{Comparison of the two Lipkin mappings}
\begin{center}
\begin{tabular}{lcc}
\hline
Case & Qubits & Main feature \\
\hline
\(N=4\), compressed & 3 & minimal encoding \\
\(N=6\), direct & 6 & simplest Pauli structure \\
\hline
\end{tabular}
\end{center}

Compressed mappings save qubits but require code-space control.
Direct mappings use more qubits but give very transparent Hamiltonians.
\end{frame}

% =========================================================
\section{Part II: Fermions, Jordan--Wigner, and Pauli matrices}
% =========================================================

\begin{frame}{Second quantization for fermions}
In fermionic many-body theory, the basic operators are
\[
a_p^\dagger,\qquad a_p,
\]
satisfying anticommutation relations
\[
\{a_p,a_q^\dagger\}=\delta_{pq},
\qquad
\{a_p,a_q\}=0,
\qquad
\{a_p^\dagger,a_q^\dagger\}=0.
\]

A general second-quantized Hamiltonian has the form
\[
H
=
\sum_{pq} h_{pq} a_p^\dagger a_q
+
\frac{1}{4}\sum_{pqrs} v_{pqrs} a_p^\dagger a_q^\dagger a_s a_r.
\]
\end{frame}

\begin{frame}{Occupation-number basis and qubits}
A fermionic Fock basis state is specified by occupations
\[
\ket{n_0 n_1 \cdots n_{M-1}},
\qquad
n_p\in\{0,1\}.
\]

This is already a binary basis, so one may map it directly to \(M\) qubits:
\[
\ket{n_0 n_1 \cdots n_{M-1}}_{\mathrm{Fock}}
\leftrightarrow
\ket{n_0 n_1 \cdots n_{M-1}}_{\mathrm{qubit}}.
\]

The nontrivial part is not the basis mapping but the operator mapping:
fermionic signs must be reproduced correctly on qubits.
\end{frame}

\begin{frame}{Why fermion-to-qubit mappings are needed}
A quantum computer acts on qubits, not directly on fermionic Fock states.
So we need a mapping
\[
a_p^\dagger,\; a_p,\; n_p=a_p^\dagger a_p
\quad \longrightarrow \quad
\text{Pauli strings}.
\]

The most common constructions are:
\begin{itemize}
\item Jordan--Wigner,
\item Bravyi--Kitaev,
\item parity mapping.
\end{itemize}

We focus here on Jordan--Wigner, since it is the most direct mathematically.
\end{frame}

\begin{frame}{Deriving the Jordan--Wigner idea}
For a single qubit,
\[
\sigma^+ = \frac{X-iY}{2} = \ket{1}\bra{0},
\qquad
\sigma^- = \frac{X+iY}{2} = \ket{0}\bra{1}.
\]

If fermions were distinguishable, we could simply identify
\[
a_p^\dagger \sim \sigma_p^+,\qquad a_p \sim \sigma_p^-.
\]

But this fails because different fermionic modes must anticommute:
\[
a_p a_q = - a_q a_p \quad (p\neq q).
\]

So we must insert a parity string that counts the number of occupied orbitals before site \(p\).
\end{frame}

\begin{frame}{Jordan--Wigner transformation}
For orbital \(p\), the Jordan--Wigner mapping is
\[
a_p^\dagger
=
\left(\prod_{j=0}^{p-1} Z_j\right)\frac{X_p-iY_p}{2},
\qquad
a_p
=
\left(\prod_{j=0}^{p-1} Z_j\right)\frac{X_p+iY_p}{2}.
\]

The string of \(Z\)'s ensures the correct fermionic sign structure.
It is often called the \emph{parity string}.
\end{frame}

\begin{frame}{Why the parity string works}
Recall that
\[
Z \ket{0} = +\ket{0},\qquad Z\ket{1} = -\ket{1}.
\]

Therefore
\[
\prod_{j=0}^{p-1} Z_j
\]
returns a factor
\[
(-1)^{n_0+n_1+\cdots+n_{p-1}}
\]
on a basis state.

This is exactly the fermionic sign picked up when a creation or annihilation operator is moved through all earlier occupied orbitals in the canonical ordering.
\end{frame}

\begin{frame}{Checking the canonical anticommutator}
Using the Jordan--Wigner map, one can verify
\[
a_p a_p^\dagger + a_p^\dagger a_p = I.
\]

Indeed, since the parity string squares to the identity,
\[
a_p^\dagger a_p
=
\frac{1-Z_p}{2},
\qquad
a_p a_p^\dagger
=
\frac{1+Z_p}{2},
\]
and therefore
\[
a_p a_p^\dagger + a_p^\dagger a_p
=
\frac{1+Z_p}{2}+\frac{1-Z_p}{2}=I.
\]
\end{frame}

\begin{frame}{Number operators under Jordan--Wigner}
A key identity is
\[
n_p = a_p^\dagger a_p = \frac{1-Z_p}{2}.
\]

Thus one-body diagonal terms become sums of \(I\) and \(Z\):
\[
\sum_p \epsilon_p a_p^\dagger a_p
\quad \longrightarrow \quad
\sum_p \epsilon_p \frac{1-Z_p}{2}.
\]

This is one of the most useful simplifications in practical qubit Hamiltonians.
\end{frame}

\begin{frame}{Nearest-neighbor hopping example}
For \(p<q\), one finds
\[
a_p^\dagger a_q + a_q^\dagger a_p
\]
maps to
\[
\frac12
\left(
X_p Z_{p+1}\cdots Z_{q-1} X_q
+
Y_p Z_{p+1}\cdots Z_{q-1} Y_q
\right).
\]

So hopping terms become nonlocal Pauli strings because of the fermionic parity chain.
\end{frame}

\begin{frame}{Pair creation operator in full orbital space}
Consider a time-reversed pair \((p,\bar p)\) with \(p<\bar p\).
Then
\[
P_p^\dagger = a_p^\dagger a_{\bar p}^\dagger.
\]

Using Jordan--Wigner,
\[
a_p^\dagger a_{\bar p}^\dagger
=
\left(\prod_{j=0}^{p-1} Z_j\right)\frac{X_p-iY_p}{2}
\left(\prod_{j=0}^{\bar p-1} Z_j\right)\frac{X_{\bar p}-iY_{\bar p}}{2}.
\]

After cancelling the common prefix, one gets a Pauli string with a parity chain between \(p\) and \(\bar p\).
\end{frame}

\begin{frame}{Explicit formula for \(a_p^\dagger a_q^\dagger\)}
For \(p<q\),
\[
a_p^\dagger a_q^\dagger
=
\frac14
(X_p-iY_p)\,
Z_{p+1}\cdots Z_{q-1}\,
(X_q-iY_q).
\]

Expanding this product gives
\[
a_p^\dagger a_q^\dagger
=
\frac14\Big(
X_p Z\cdots Z X_q
- i X_p Z\cdots Z Y_q
- i Y_p Z\cdots Z X_q
- Y_p Z\cdots Z Y_q
\Big),
\]
where the intermediate \(Z\)-string runs from \(p+1\) to \(q-1\).
\end{frame}

\begin{frame}{Hermitian pair transfer operator}
A very common Hermitian combination is
\[
a_p^\dagger a_q^\dagger + a_q a_p.
\]

Under Jordan--Wigner this becomes
\[
\frac12
\left(
X_p Z_{p+1}\cdots Z_{q-1} X_q
-
Y_p Z_{p+1}\cdots Z_{q-1} Y_q
\right).
\]

This is the pairing analogue of the familiar hopping identity, with a sign change between the \(XX\) and \(YY\) channels.
\end{frame}

% =========================================================
\section{Part III: Quasispin algebra from second quantization}
% =========================================================

\begin{frame}{Local pairing operators}
For each time-reversed pair \((p,\bar p)\), define
\[
P_p^\dagger = a_p^\dagger a_{\bar p}^\dagger,
\qquad
P_p = a_{\bar p} a_p,
\]
and the pair number operator
\[
N_p = a_p^\dagger a_p + a_{\bar p}^\dagger a_{\bar p}.
\]

Now define local quasispin operators
\[
S_p^+ = P_p^\dagger,
\qquad
S_p^- = P_p,
\qquad
S_p^z = \frac12(N_p-1).
\]
\end{frame}

\begin{frame}{Commutator \([S_p^z,S_p^+]\)}
We derive the quasispin algebra directly from fermion operators.

Using
\[
S_p^z = \frac12(N_p-1),\qquad S_p^+ = a_p^\dagger a_{\bar p}^\dagger,
\]
we compute
\[
[S_p^z,S_p^+]
=
\frac12[N_p,a_p^\dagger a_{\bar p}^\dagger].
\]

Since
\[
[N_p,a_p^\dagger]=a_p^\dagger,\qquad [N_p,a_{\bar p}^\dagger]=a_{\bar p}^\dagger,
\]
it follows that
\[
[N_p,a_p^\dagger a_{\bar p}^\dagger]
=
2 a_p^\dagger a_{\bar p}^\dagger.
\]

Hence
\[
[S_p^z,S_p^+] = S_p^+.
\]
\end{frame}

\begin{frame}{Commutator \([S_p^z,S_p^-]\)}
Similarly,
\[
[S_p^z,S_p^-]
=
\frac12[N_p,a_{\bar p}a_p].
\]

Using
\[
[N_p,a_p]=-a_p,\qquad [N_p,a_{\bar p}]=-a_{\bar p},
\]
we find
\[
[N_p,a_{\bar p}a_p] = -2 a_{\bar p}a_p.
\]

Therefore
\[
[S_p^z,S_p^-] = -S_p^-.
\]
\end{frame}

\begin{frame}{Commutator \([S_p^+,S_p^-]\)}
Now compute
\[
[S_p^+,S_p^-]
=
a_p^\dagger a_{\bar p}^\dagger a_{\bar p} a_p
-
a_{\bar p} a_p a_p^\dagger a_{\bar p}^\dagger.
\]

The first term is
\[
a_p^\dagger a_{\bar p}^\dagger a_{\bar p} a_p
=
a_p^\dagger n_{\bar p} a_p
=
n_p n_{\bar p}.
\]

The second term becomes
\[
a_{\bar p}(1-n_p)a_{\bar p}^\dagger
=
(1-n_p)(1-n_{\bar p}).
\]

Subtracting yields
\[
[S_p^+,S_p^-]
=
n_p+n_{\bar p}-1
=
N_p-1
=
2S_p^z.
\]
\end{frame}

\begin{frame}{Local quasispin algebra}
We have proved
\[
[S_p^z,S_p^\pm] = \pm S_p^\pm,
\qquad
[S_p^+,S_p^-] = 2S_p^z.
\]

So for each pair level \(p\), the operators
\[
S_p^+,\qquad S_p^-,\qquad S_p^z
\]
generate an \(su(2)\) algebra.

For different levels \(p\neq q\), the operators commute:
\[
[S_p^\alpha,S_q^\beta]=0.
\]

Thus a pairing model is naturally a sum of commuting local quasispins.
\end{frame}

\begin{frame}{Collective quasispin}
Defining
\[
S^\pm = \sum_p S_p^\pm,
\qquad
S^z = \sum_p S_p^z,
\]
we immediately get the collective algebra
\[
[S^z,S^\pm]=\pm S^\pm,
\qquad
[S^+,S^-]=2S^z.
\]

This is the precise sense in which the pairing problem can be rewritten as a spin problem.
It also explains the close analogy with the Lipkin Hamiltonian.
\end{frame}

% =========================================================
\section{Part IV: Nuclear pairing Hamiltonian}
% =========================================================

\begin{frame}{Nuclear pairing Hamiltonian}
A standard reduced pairing Hamiltonian is
\[
H_{\mathrm{pair}}
=
\sum_{p} \epsilon_p \, N_p
-
G \sum_{p,q} P_p^\dagger P_q,
\]
where
\[
N_p = a_p^\dagger a_p + a_{\bar p}^\dagger a_{\bar p},
\]
\[
P_p^\dagger = a_p^\dagger a_{\bar p}^\dagger,
\qquad
P_p = a_{\bar p} a_p.
\]

Here \(p\) and \(\bar p\) denote time-reversed partners.
\end{frame}

\begin{frame}{Pairing Hamiltonian in quasispin form}
Using
\[
S_p^+ = P_p^\dagger,\qquad
S_p^- = P_p,\qquad
S_p^z = \frac12(N_p-1),
\]
the pairing Hamiltonian becomes
\[
H_{\mathrm{pair}}
=
\sum_p \epsilon_p (2S_p^z+1)
-
G \sum_{p,q} S_p^+ S_q^-.
\]

Thus the pairing Hamiltonian is an interacting quasispin Hamiltonian.
This is the natural starting point for compact qubit encodings.
\end{frame}

\begin{frame}{One pair level as one qubit}
For a single pair \((p,\bar p)\), the paired subspace is two-dimensional:
\[
\ket{0}_p \equiv \text{pair unoccupied},
\qquad
\ket{1}_p \equiv \text{pair occupied}.
\]

In this reduced subspace, we may identify
\[
S_p^+ = \frac{X_p-iY_p}{2},
\qquad
S_p^- = \frac{X_p+iY_p}{2},
\qquad
S_p^z = -\frac{Z_p}{2}.
\]

Also,
\[
N_p = 1-Z_p.
\]
\end{frame}

\begin{frame}{Pairing Hamiltonian in pair-qubit form}
Substituting
\[
N_p = 1-Z_p,
\qquad
P_p^\dagger = \frac{X_p-iY_p}{2},
\qquad
P_p = \frac{X_p+iY_p}{2},
\]
we obtain
\[
H_{\mathrm{pair}}
=
\sum_p \epsilon_p(1-Z_p)
-
G\sum_{p,q} S_p^+S_q^-.
\]

Since
\[
S_p^+S_q^- + S_q^+S_p^-
=
\frac12(X_pX_q+Y_pY_q),
\]
the interaction becomes an \(XY\)-type coupling between pair qubits.
\end{frame}

\begin{frame}{Explicit pairing Hamiltonian for \(\Omega\) pair levels}
For \(\Omega\) pair orbitals, the pairing Hamiltonian in the paired subspace can be written as
\[
H_{\mathrm{pair}}
=
\sum_{p=1}^{\Omega} \epsilon_p (1-Z_p)
-
\frac{G}{2}\sum_{p\ne q}(X_pX_q+Y_pY_q)
-
G\sum_p S_p^+S_p^-.
\]

Since
\[
S_p^+S_p^- = \frac{1-Z_p}{2},
\]
this yields
\[
H_{\mathrm{pair}}
=
\sum_{p=1}^{\Omega}\left(\epsilon_p-\frac{G}{2}\right)(1-Z_p)
-
\frac{G}{2}\sum_{p\ne q}(X_pX_q+Y_pY_q).
\]
\end{frame}

\begin{frame}{Example: two pair levels in pair-qubit space}
For two pair orbitals \(p=1,2\), the paired-subspace Hamiltonian becomes
\[
H_{\mathrm{pair}}^{(2)}
=
\left(\epsilon_1-\frac{G}{2}\right)(1-Z_1)
+
\left(\epsilon_2-\frac{G}{2}\right)(1-Z_2)
-
\frac{G}{2}(X_1X_2+Y_1Y_2).
\]

This is a compact two-qubit Hamiltonian ready for variational or time-evolution algorithms.
\end{frame}

\begin{frame}{Full orbital-space example: one pair level}
Now keep the two fermionic orbitals \(p\) and \(\bar p\) explicitly as two qubits.
The one-level pairing Hamiltonian is
\[
H = \epsilon_p(n_p+n_{\bar p}) - G\, a_p^\dagger a_{\bar p}^\dagger a_{\bar p} a_p.
\]

Using Jordan--Wigner,
\[
n_p=\frac{1-Z_p}{2},\qquad
n_{\bar p}=\frac{1-Z_{\bar p}}{2}.
\]

Also
\[
a_p^\dagger a_{\bar p}^\dagger a_{\bar p} a_p = n_p n_{\bar p}.
\]

Therefore
\[
H
=
\frac{\epsilon_p}{2}(2 - Z_p - Z_{\bar p})
-
\frac{G}{4}(1-Z_p)(1-Z_{\bar p}).
\]
\end{frame}

\begin{frame}{Expanded one-pair full-orbital Pauli form}
Expanding the previous result gives
\[
H
=
\left(\epsilon_p-\frac{G}{4}\right)I
+
\left(-\frac{\epsilon_p}{2}+\frac{G}{4}\right)Z_p
+
\left(-\frac{\epsilon_p}{2}+\frac{G}{4}\right)Z_{\bar p}
-
\frac{G}{4} Z_p Z_{\bar p}.
\]

So even the full-orbital one-pair Hamiltonian reduces to a simple two-qubit diagonal Pauli Hamiltonian.
\end{frame}

\begin{frame}{Full orbital-space example: pair transfer between two levels}
Consider two pair levels \(p\) and \(q\), with orbitals
\[
(p,\bar p,q,\bar q).
\]

A pairing transfer term is
\[
P_p^\dagger P_q + P_q^\dagger P_p
=
a_p^\dagger a_{\bar p}^\dagger a_{\bar q} a_q
+
a_q^\dagger a_{\bar q}^\dagger a_{\bar p} a_p.
\]

Under Jordan--Wigner, this becomes a Hermitian sum of four-qubit Pauli strings.
\end{frame}

\begin{frame}{Explicit Pauli decomposition of pair transfer}
Assume the orbital ordering
\[
p < \bar p < q < \bar q.
\]

Then one finds
\[
P_p^\dagger P_q + P_q^\dagger P_p
=
\frac{1}{8}
\Big(
X_p X_{\bar p} X_q X_{\bar q}
-
X_p Y_{\bar p} Y_q X_{\bar q}
\]
\[
\qquad
-
Y_p X_{\bar p} X_q Y_{\bar q}
+
Y_p Y_{\bar p} Y_q Y_{\bar q}
\Big)
\]
up to the parity-string conventions induced by the chosen ordering.
In adjacent pair ordering, the intermediate \(Z\)-chains cancel within the paired sector, leaving a compact four-body expression.
\end{frame}

\begin{frame}{Interpretation of the full-orbital pairing terms}
This explicit decomposition shows:
\begin{itemize}
\item full fermionic orbital space leads to longer Pauli strings,
\item pair-qubit reduction produces shorter and more structured qubit Hamiltonians,
\item both are exact, but they operate in different Hilbert spaces.
\end{itemize}

So there is a practical tradeoff:
\begin{itemize}
\item full-space mappings are fully general,
\item reduced quasispin/pair encodings are more compact when the physics is dominated by seniority-zero pairing.
\end{itemize}
\end{frame}

\begin{frame}{Lipkin versus pairing: structural comparison}
Both Hamiltonians admit spin-like descriptions:
\begin{itemize}
\item Lipkin:
\[
H_{\mathrm{Lipkin}}
=
\frac{\epsilon}{2}\left(J_z-\frac{N}{2}\right)
+\frac{V}{2}(J_+^2+J_-^2)
+\frac{W}{2}J_z^2,
\]
\item Pairing:
\[
H_{\mathrm{pair}}
=
\sum_p \epsilon_p N_p - G\sum_{p,q} P_p^\dagger P_q.
\]
\end{itemize}

Both become sums of Pauli strings after:
\begin{itemize}
\item choosing an encoding,
\item expressing collective or pair operators in terms of \(X,Y,Z\).
\end{itemize}
\end{frame}

\begin{frame}{General recipe: from second quantization to qubits}
For a fermionic Hamiltonian:
\begin{enumerate}
\item write the Hamiltonian in terms of \(a_p^\dagger,a_p\),
\item decide whether a reduced quasispin or pair-qubit subspace is appropriate,
\item otherwise apply Jordan--Wigner or a related mapping,
\item simplify number and pair operators,
\item collect the resulting Pauli strings.
\end{enumerate}

For collective-spin models:
\begin{enumerate}
\item write the model in \(J_x,J_y,J_z\) or ladder-operator form,
\item choose compressed or direct qubit encoding,
\item expand diagonal and off-diagonal matrix elements into Pauli strings.
\end{enumerate}
\end{frame}

\begin{frame}{Summary}
In this expanded version we added:
\begin{itemize}
\item a fuller Jordan--Wigner derivation,
\item a formal derivation of the quasispin algebra directly from second quantization,
\item explicit Pauli decompositions for small pairing examples in full orbital space,
\item and a more systematic comparison of compact versus full-space encodings.
\end{itemize}

The unifying message is:
\[
\text{fermions or collective spins}
\;\longrightarrow\;
\text{spin algebra}
\;\longrightarrow\;
\text{Pauli strings}.
\]
\end{frame}

\begin{frame}[plain,fragile]
\frametitle{Quantum Fourier transform: Overarching motivation}

We will start with a reminder on Fourier theory and in particular on discrete Fourier transforms (DFTs).  There are many
motivations for the DFTs. For those of you familiar with signal
processing, harmonic oscillations, and many other areas of
applications, Fourier transforms are almost standard kitchen
items. 

But first some general motivation for Quantum Fourier Transforms (QFTs).
\end{frame}

\begin{frame}[plain,fragile]
\frametitle{Quantum Fourier Transforms (QFTs)}

\begin{enumerate}
\item QFTs are the quantum analogue of the Discrete Fourier Transforms (DFTs).

\item They  play a crucial role in quantum algorithms like Shor’s Algorithm and Quantum Phase Estimation.

\item QFTs provide an \emph{exponential speedup} over classical Fourier Transform methods.
\end{enumerate}

\noindent
\end{frame}

\begin{frame}[plain,fragile]
\frametitle{Why Quantum Fourier Transforms and exponential speedup}

\begin{enumerate}
\item Classical Discrete Fourier Transforms (DFTs) require $O(N^2)$ operations.

\item The Fast Fourier Transform (FFT) improves this to $O(N \log N)$ operations.

\item QFTs reduce complexity to $O((\log N)^2)$ using quantum gates.
\end{enumerate}

\noindent
\end{frame}

\begin{frame}[plain,fragile]
\frametitle{Quantum Fourier Transforms and quantum parallelism}

\begin{enumerate}
\item QFTs act on a \emph{superposition} of states, processing all inputs simultaneously.

\item This is crucial in algorithms like:
\begin{enumerate}

 \item Shor’s Algorithm for factoring large numbers.

 \item Quantum Phase Estimation (QPE) for eigenvalue extraction.
\end{enumerate}

\noindent
\end{enumerate}

\noindent
\end{frame}

\begin{frame}[plain,fragile]
\frametitle{Quantum Fourier Transforms and implementations}

\begin{enumerate}
\item QFTs require only \emph{Hadamard gates and controlled-phase gates}.

\item A 3-qubit QFT circuit uses only $O(n^2)$.

\item This makes QFTs highly efficient for \emph{quantum hardware}.
\end{enumerate}

\noindent
\end{frame}

\begin{frame}[plain,fragile]
\frametitle{Quantum Fourier Transform and Applications in Quantum Computing}

\begin{enumerate}
\item \textbf{Shor’s Algorithm:} Uses QFTs to find periodicity in modular exponentiation.

\item \textbf{Quantum Phase Estimation (QPE):} QFTs extract eigenvalues of unitary matrices.

\item \textbf{Quantum Signal Processing:} QFTs enable spectral analysis on quantum data. Important for AI applications.
\end{enumerate}

\noindent
\end{frame}

\begin{frame}[plain,fragile]
\frametitle{Quantum Fourier Transforms and Reduced Memory Complexity (not covered)}

\begin{enumerate}
\item Classical DFTs require $O(N)$ space.

\item QFT store Fourier-transformed coefficients \emph{implicitly} in qubit amplitudes.

\item Only $O(n)$ qubits are needed for a size $N = 2^n$ transformation.
\end{enumerate}

\noindent
\end{frame}

\begin{frame}[plain,fragile]
\frametitle{Quantum Fourier Transforms and Quantum Signal Processing (not covered in this course)}

QFTs enable applications such as:
\begin{enumerate}
\item Quantum spectral analysis.

\item Quantum image processing.

\item Quantum filtering and denoising.
\end{enumerate}

\noindent
These can enhance AI, cryptography, and data processing.
\end{frame}

\begin{frame}[plain,fragile]
\frametitle{Why Quantum Fourier Transforms? Brief summary}

\begin{enumerate}
\item Quantum Fourier Transforms provide \emph{exponential speedup}

\item They enable key quantum algorithms like \emph{Shor’s Algorithm} and \emph{Quantum Phase Estimation}.

\item QFTs are  more efficient in \emph{memory usage and circuit complexity} than classical methods.

\item Future quantum applications in \emph{signal processing, AI, and cryptography} will heavily rely on QFT.
\end{enumerate}

\noindent
\end{frame}

\end{document}
